% ============================================================
\chapter*{Pr\'eface}
\addcontentsline{toc}{chapter}{Pr\'eface}
% ============================================================

Ce cours de \emph{Topologie Ponctuelle Avanc\'ee} s'adresse aux \'etudiants de Master~2 et
de doctorat en math\'ematiques pures, ainsi qu'aux chercheurs en informatique th\'eorique
souhaitant approfondir les fondements topologiques de la th\'eorie des domaines.

\section*{Motivations}

La topologie g\'en\'erale, telle qu'elle est enseign\'ee dans les cursus de licence et de
premi\`ere ann\'ee de master, se concentre presque exclusivement sur les espaces de Hausdorff
(s\'epar\'es). Cette restriction, bien que naturelle pour l'analyse et la g\'eom\'etrie
diff\'erentielle, occulte une partie consid\'erable de la th\'eorie. En effet, de nombreuses
structures fondamentales en math\'ematiques et en informatique ne satisfont pas l'axiome de
s\'eparation~$T_2$~:

\begin{itemize}
  \item Le spectre d'un anneau commutatif, muni de la topologie de Zariski, est
    g\'en\'eralement un espace~$T_0$ non~$T_1$.
  \item Les domaines de Scott, essentiels en s\'emantique d\'enotationnelle, portent une
    topologie~$T_0$ non~$T_1$.
  \item Les locales, qui g\'en\'eralisent la notion d'espace topologique en \'eliminant la
    r\'ef\'erence aux points, constituent un cadre naturel pour la logique constructive et la
    th\'eorie des topos.
  \item Les quasi-m\'etriques, o\`u la sym\'etrie est abandonn\'ee, mod\'elisent des notions de
    co\^ut asym\'etrique en informatique et en optimisation.
\end{itemize}

L'objectif de ce cours est de d\'evelopper syst\'ematiquement la th\'eorie topologique dans ce
cadre \'elargi, en mettant l'accent sur les interactions entre topologie, th\'eorie de l'ordre
et informatique th\'eorique.

\section*{Organisation du cours}

Le cours est divis\'e en dix chapitres, organis\'es en trois parties th\'ematiques~:

\medskip
\noindent\textbf{Partie~I~: Fondements (Chapitres~1--3).}
Nous commen\c{c}ons par un rappel syst\'ematique de topologie g\'en\'erale enrichi du langage
cat\'egoriel (Chapitre~1), puis nous \'etudions les espaces~$T_0$ via l'ordre de
sp\'ecialisation, \'etablissant le pont fondamental entre topologie et th\'eorie de l'ordre
(Chapitre~2), avant d'aborder les espaces sobres et la th\'eorie des locales qui fournissent
le cadre <<~sans points~>> de la topologie (Chapitre~3).

\medskip
\noindent\textbf{Partie~II~: Structures asym\'etriques (Chapitres~4--6).}
Nous d\'eveloppons la th\'eorie des quasi-m\'etriques et de la topologie asym\'etrique
(Chapitre~4), des espaces bitopologiques qui capturent la dualit\'e entre une topologie et
sa conjugu\'ee (Chapitre~5), puis la topologie de Scott sur les ensembles ordonn\'es et les
domaines, fondamentale pour la s\'emantique d\'enotationnelle (Chapitre~6).

\medskip
\noindent\textbf{Partie~III~: Convergence, compacit\'e et applications (Chapitres~7--10).}
Nous \'etudions les espaces de convergence g\'en\'eralis\'es et la th\'eorie des filtres
(Chapitre~7), la compacit\'e en dehors du cadre de Hausdorff avec ses multiples variantes
(Chapitre~8), les applications de la topologie asym\'etrique \`a l'informatique th\'eorique
(Chapitre~9), et enfin la th\'eorie des points fixes topologiques avec ses applications aux
domaines et \`a la s\'emantique (Chapitre~10).

\section*{Pr\'erequis}

Le lecteur est suppos\'e conna\^itre~:
\begin{itemize}
  \item La topologie g\'en\'erale de base~: espaces topologiques, continuit\'e, compacit\'e,
    connexit\'e, axiomes de s\'eparation $T_0$--$T_4$ (niveau L3/M1).
  \item Les notions \'el\'ementaires de th\'eorie des ensembles ordonn\'es~: ordres partiels,
    treillis, compl\'etude.
  \item Les bases du langage cat\'egoriel~: cat\'egories, foncteurs, transformations
    naturelles, adjonctions.
  \item La th\'eorie \'el\'ementaire des espaces m\'etriques et de la convergence.
\end{itemize}

\section*{Conventions et notations}

Tout au long de ce cours, nous adoptons les conventions suivantes~:
\begin{itemize}
  \item $(X,\tau)$ d\'esigne un espace topologique, $\tau$ \'etant la topologie (ensemble des
    ouverts). Lorsque la topologie est claire du contexte, nous \'ecrivons simplement~$X$.
  \item $\mathcal{O}(X)$ d\'esigne le treillis des ouverts de~$X$.
  \item $\mathcal{C}(X)$ d\'esigne l'ensemble des ferm\'es de~$X$.
  \item $\overline{A}$, $\mathrm{Int}(A)$ d\'esignent respectivement l'adh\'erence et
    l'int\'erieur de $A \subseteq X$.
  \item $\leq_X$ ou $\sqsubseteq$ d\'esigne l'ordre de sp\'ecialisation sur un espace~$T_0$.
  \item $\downarrow\! x = \{y : y \leq x\}$ et $\uparrow\! x = \{y : x \leq y\}$.
  \item Pour un sous-ensemble $A$, on note $\downarrow\! A = \bigcup_{a\in A}\downarrow\! a$
    et $\uparrow\! A = \bigcup_{a\in A}\uparrow\! a$.
  \item $\mathbf{Top}$ est la cat\'egorie des espaces topologiques et applications continues.
  \item $\mathbf{Top}_0$ est la sous-cat\'egorie pleine des espaces~$T_0$.
  \item $\mathbf{Sob}$ est la sous-cat\'egorie pleine des espaces sobres.
  \item $\mathbf{Loc}$ est la cat\'egorie des locales.
  \item $\mathbf{Frm}$ est la cat\'egorie des cadres (\emph{frames}).
  \item $\mathbf{Dcpo}$ est la cat\'egorie des dcpo et fonctions Scott-continues.
  \item $\mathbf{Dom}$ d\'esigne la cat\'egorie des domaines (selon le contexte).
  \item $\omega$ d\'esigne le premier ordinal infini, identifi\'e \`a $\N$.
  \item $\mathbb{S}$ d\'esigne l'espace de Sierpi\'nski $\{0,1\}$ avec la topologie
    $\{\emptyset, \{1\}, \{0,1\}\}$.
\end{itemize}

\section*{Approche p\'edagogique}

Chaque chapitre contient~:
\begin{itemize}
  \item Des \textbf{d\'efinitions} pr\'ecises et des \textbf{th\'eor\`emes} avec d\'emonstrations
    compl\`etes.
  \item Des \textbf{exemples} d\'etaill\'es et des \textbf{contre-exemples} illustrant les
    limites des r\'esultats.
  \item Des \textbf{exercices} de difficult\'e vari\'ee, certains prolong\'es par des indications.
  \item Des encadr\'es \textbf{Intuition} et \textbf{Attention} pour guider la compr\'ehension.
\end{itemize}

\section*{Sources et r\'ef\'erences}

Ce cours s'appuie principalement sur les ouvrages suivants~:
\begin{itemize}
  \item \textsc{Gierz, Hofmann, Keimel, Lawson, Mislove, Scott},
    \emph{Continuous Lattices and Domains}, Cambridge University Press, 2003.
  \item \textsc{Goubault-Larrecq}, \emph{Non-Hausdorff Topology and Domain Theory},
    Cambridge University Press, 2013.
  \item \textsc{Johnstone}, \emph{Stone Spaces}, Cambridge University Press, 1982.
  \item \textsc{Vickers}, \emph{Topology via Logic}, Cambridge University Press, 1989.
  \item \textsc{Abramsky, Jung}, \emph{Domain Theory}, Handbook of Logic in Computer
    Science, vol.~3, Oxford University Press, 1994.
  \item \textsc{Kelly}, \emph{Bitopological Spaces}, Proc.\ London Math.\ Soc., 1963.
  \item \textsc{Fletcher, Lindgren}, \emph{Quasi-Uniform Spaces}, Marcel Dekker, 1982.
  \item \textsc{Smyth}, \emph{Topology}, Handbook of Logic in Computer Science, vol.~1,
    Oxford University Press, 1992.
  \item \textsc{Escard\'o}, \emph{Synthetic Topology of Data Types and Classical Spaces},
    ENTCS, 2004.
  \item \textsc{Nachbin}, \emph{Topology and Order}, Van Nostrand, 1965.
\end{itemize}

\section*{Remerciements}

Ce cours a b\'en\'efici\'e des enseignements et des discussions avec de nombreux coll\`egues
travaillant \`a l'interface de la topologie, de la th\'eorie de l'ordre et de l'informatique
th\'eorique. Nous remercions en particulier les \'etudiants dont les questions et remarques
ont contribu\'e \`a am\'eliorer la pr\'esentation.

\bigskip
\hfill\emph{L'auteur}
